\documentclass[11pt]{article}
\usepackage{graphicx}
\usepackage{amssymb}
\usepackage{bm}
\usepackage{mathtools}
\usepackage{amsmath}
\usepackage{commath}
\usepackage{amsthm}
\usepackage{enumitem}
\usepackage{float} 
\usepackage{array}
\usepackage{outlines}
\usepackage{setspace}
\usepackage[colorinlistoftodos]{todonotes}
\definecolor{green}{rgb}{0.4, 0.69, 0.2}
\definecolor{asparagus}{rgb}{0.53, 0.66, 0.42}
\definecolor{blue}{rgb}{0.0, 0.53, 0.74}
\setlist[enumerate,1]{label=\roman*)}
\setlist[enumerate,2]{label=\alph*)}
\newtheorem{theorem}{Theorem}
\newtheorem{corollary}{Corollary}
\newtheorem{lemma}{Lemma}
\newtheorem{remark}{Remark}
\newtheorem{definition}{Definition}
\newtheorem{prop}{Proposition}
\newcommand*\diff{\mathop{}\!\mathrm{d}}
\usepackage{tikz}
\usetikzlibrary{arrows,shapes,trees, calc}


%%% *** PREAMBLE *** %%%
\title{Strong solutions}

\date{}

%%%***  CONTENT *** %%%
\begin{document}

\maketitle

\begin{abstract}

The strong convergence of the strong solutions.

\end{abstract}

\newpage
\tableofcontents
\newpage

\section{The case of non-periodic boundary for the strong solutions}

See the calculations on the paper, it seems that we need a periodic boundary.

\section{Geophysical coordinates}

\subsection{The collection of the equations we work with and some definitions, results we have}

Firstly, here we work with the geophysical coordinate system and $K_1, K_2, K_3$ diffusion terms.

\textbf{Change of notations:} Since $u^\epsilon$ and $u$ denotes the solution functions and the capital $U^\epsilon$ denotes their difference, we switch to the analogous notation with the concentration function, i.e. $c^\epsilon$ and $c$ here are the solutions in the anisotropic and hydrostatic equations, respectively, while the capital $C^\epsilon$ denotes the $c^\epsilon - c$ difference function.

(E-SNS): extended scaled Navier-Stokes.

    \begin{equation} \label{ANSC1}
    \partial_t u^\epsilon _1 + \mathbf{u}^\epsilon \cdot \nabla u^\epsilon _1 -\Delta_\nu u^\epsilon _1 -\gamma u_2^\epsilon + \epsilon \beta u_3^\epsilon + \partial_1 p^\epsilon = 0 \text{ in } \Omega \times (0,T)
    \end{equation}
    
    \begin{equation} \label{ANSC2}
    \partial_t u_2^\epsilon + \mathbf{u}^\epsilon \cdot \nabla u_2^\epsilon -\Delta_\nu u_2^\epsilon + \gamma u_1^\epsilon   + \partial_2 p^\epsilon = 0 \text{ in } \Omega \times (0,T)
    \end{equation}
    
    \begin{equation} \label{ANSC3}
    \epsilon^2 \{  \partial_t u_3^\epsilon + \mathbf{u}^\epsilon \cdot \nabla u_3^\epsilon -\Delta_\nu u_3^{\epsilon} \} - \epsilon \beta u_1^\epsilon + \partial_3 p^\epsilon = 0 \text{ in } \Omega \times (0,T)
    \end{equation}
    
    \begin{equation} \label{ANSC4}
    \nabla \cdot \mathbf{u}^\epsilon= 0 \text{ in } \Omega \times (0,T)
    \end{equation}
    
    \begin{equation} \label{ANSC5}
    \partial_t c^\epsilon + \mathbf{u^\epsilon} \cdot \nabla c^\epsilon = K_1 \partial_{11} ^ 2 c^\epsilon + K_2 \partial_{22} ^ 2 c^\epsilon + K_3 \partial_{33} ^ 2 c^\epsilon + s^\epsilon \text{ in } \Omega \times (0,T).
    \end{equation}

(E-PEs): extended Primitive Equations.

    \begin{equation} \label{HNSC1}    
    \partial_t u_1 + \mathbf{u} \cdot \nabla u_1 - \Delta_\nu u_1 - \gamma u_2 + \partial_1 p = 0 \text{ in } \Omega \times (0,T)
    \end{equation}
    
    \begin{equation} \label{HNSC2}   
    \partial_t u_2 + \mathbf{u} \cdot \nabla u_2 - \Delta_\nu u_2 + \gamma u_1 + \partial_2 p = 0 \text{ in } \Omega \times (0,T)
    \end{equation}
    
    \begin{equation} \label{HNSC3}   
    \partial_3 p = 0 \text{ in } \Omega \times (0,T)
    \end{equation}
    
    \begin{equation} \label{HNSC4}   
    \nabla \cdot \mathbf{u} = 0 \text{ in } \Omega \times (0,T)
    \end{equation}
    
    \begin{equation} \label{HNSC5}   
    \partial_t c + \mathbf{u} \cdot \nabla c = K_1 \partial_{11} ^ 2 c + K_2 \partial_{22} ^ 2 c + K_3 \partial_{33} ^2 c + s \text{ in } \Omega \times (0,T)
    \end{equation}

The energy inequality for the geophysical coordinate system for a periodic boundary becomes more simple (originally we used the Trace Theorem to handle the boundary term, here we do not need that.) We use $\mathbf{K} = (K_1, K_2, K_3).$

\begin{equation}
  \begin{split}  
    & \frac{1}{2} \big (\norm{\mathbf{u}_H^\epsilon (t)}_{L^2}^2 + \epsilon ^ 2 \norm{u_3^\epsilon (t)}_{L^2}^2 + \norm{c^\epsilon (t)}_{L^2}^2 \big ) \\
    & + \int_0 ^ t \bigg [\norm{\nabla_\nu \mathbf{u}_H^\epsilon}_{L^2}^2 + \epsilon ^ 2 \norm{\nabla_\nu u_3^\epsilon}_{L^2}^2 + \norm{ \mathbf{K} \nabla c^\epsilon}_{L^2}^2\bigg ] \diff \tau \\
    & \leq \frac{1}{2} (\norm{\mathbf{u}_H^\epsilon (0)}_{L^2}^2 + \epsilon ^ 2 \norm{u_3^\epsilon (0)}_{L^2}^2 + \norm{c^\epsilon (0)}_{L^2}^2) \\
    & + \int_0^t (s^\epsilon,c^\epsilon) \diff \tau. \label{WeakSolEnergyInequality}
  \end{split}
\end{equation}  

\subsection{Recreate the $L^2$ bound for the extended system (analogous to Section 4 in Li-Titi)}

Here we need to show that $$ \int_0^t \norm{\nabla C^\epsilon}_2^2 \leq \kappa \epsilon^2,$$ because in the next section we will assume at some point, that this is smaller than some $\delta_0.$ That makes sense only if the expression is arbitrarily small.

The expression we need to bound here comes together from four results.

ONE. an identity, (4.4) in Li-Titi

TWO.

THREE. Multiplying the (E-PEs) by $(v,c)$ and integrating. This gives the extended version of what is (4.9) in Li-Titi:

\begin{equation}
  \begin{split}
    & \frac{1}{2} \big ( \norm{v(t_0)}_2^2 + \norm{c(t_0)}_2^2 \big) + \int_0^{t_0} ( \norm{\nabla v}_2^2 + K \norm{\nabla c}_2^2 ) \diff \tau \\
    & = \frac{1}{2} \big ( \norm{v_0}_2^2 + \norm{c_0}_2^2 \big) + \int_{Q_{t_0}} s \cdot c.
  \end{split}
\end{equation}

This comes using integration by parts, using $\partial_3 p = 0,$ the divergence-free condition, and that at the integration by part process the $u \nabla v_1 \cdot v_1$ kind of terms fall out, as they have both the function and its derivative, i.e. $\nabla v_1$ and $v_1$, and $\nabla \cdot u$ is zero, and this we have that something is equals to minus something, thus it is 0.

FOUR. The original energy inequality, required by definition of the weak solution, (\ref{WeakSolEnergyInequality}). (This is (4.10) in Li-Titi)

\subsection{Recreate the full strong convergence result for the extended system (analogous to Section 5 in Li-Titi)}

With initial data $v_0 \in H^2 (\Omega), c_0 \in H^1(\Omega).$ For now we assume this, and we will see later if it is enough. We have $u_0 \in H^1(\Omega)$ in this case.
  
\textbf{Starting point / Existence:}

\begin{enumerate}
  \item The existence of a local in time strong solution $(u^\epsilon, c^\epsilon)$ to (E-SNS) with maximal existence time $T^\epsilon.$ TODO: details. We know the existence of $u^\epsilon$ for (SNS).
  \item The existence of the (u,c) strong solution globally of (E-PEs). TODO: details. We have the existence of $u.$
\end{enumerate}
  
\textbf{Convergence:}  
Denote $(U^\epsilon, C^\epsilon) = (V^\epsilon, W^\epsilon, C^\epsilon) = (u^\epsilon - u, c^\epsilon - c) = (v^\epsilon - v, w^\epsilon - w, c^\epsilon - c).$

The equations describing $(U^\epsilon, C^\epsilon)$ are obtained by \textit{a)} extending equations (5.1)-(5.3) in Li-Titi by the Coriolis terms (see in blue below) and \textit{b)} calculating the equation describing the concentration difference (in green).

Calculating the difference between (\ref{ANSC5}) and (\ref{HNSC5}) we have

$$ \partial_t C^\epsilon + u^\epsilon \cdot \nabla c^\epsilon - u \cdot \nabla c = K_1 \partial_{11} C^\epsilon + K_2 \partial_{22} C^\epsilon + K_3 \partial_{33} C^\epsilon + S^\epsilon.$$

Rewriting the term $u^\epsilon \cdot \nabla c^\epsilon - u \cdot \nabla c$ we use

\begin{equation}
  \begin{split}
    & u^\epsilon \cdot \nabla c^\epsilon - u \cdot \nabla c + 0 = u^\epsilon \cdot \nabla c^\epsilon - u \cdot \nabla c - u \nabla c^\epsilon + u \nabla c^\epsilon \\
    & = u^\epsilon \cdot \nabla c^\epsilon - u \nabla c^\epsilon - u \cdot \nabla c + u \nabla c^\epsilon = U^\epsilon \nabla c^\epsilon + u \nabla C^\epsilon \\
    & = U^\epsilon \nabla (c^\epsilon - c + c) + u \nabla C^\epsilon = U^\epsilon \nabla C^\epsilon  + U^\epsilon \nabla c+ u \nabla C^\epsilon.
  \end{split}
\end{equation}

Using this, we arrive to the following set of equation that describes $(U^\epsilon, C^\epsilon)$:

\begin{equation} \label{U_eps_C_eps}
  \begin{gathered}
    \partial_t V^\epsilon + (U^\epsilon \cdot \nabla) V^\epsilon - \Delta V^\epsilon + \nabla_h P^\epsilon + (u \cdot \nabla) V^\epsilon + (U^\epsilon \cdot \nabla) v \textcolor{blue}{- \gamma V^\epsilon + \epsilon \beta W^\epsilon  + \epsilon \beta w} = 0, \\
    \nabla_h \cdot V^\epsilon + \partial_z W^\epsilon = 0, \\
    \epsilon^2 (\partial_t W^\epsilon + U^\epsilon \cdot \nabla W^\epsilon - \Delta W^\epsilon + U^\epsilon \cdot \nabla w + u \cdot \nabla W^\epsilon) \\
    + \partial_z P^\epsilon \textcolor{blue}{- \epsilon \beta V_1^\epsilon - \epsilon \beta v_1} = - \epsilon^2 (\partial_t w + u \cdot \nabla w - \Delta w), \\
     \textcolor{green}{\partial_t C^\epsilon + U^\epsilon \nabla C^\epsilon  + U^\epsilon \nabla c+ u \nabla C^\epsilon = K_1 \partial_{11} C^\epsilon + K_2 \partial_{22} C^\epsilon + K_3 \partial_{33} C^\epsilon + S^\epsilon}
  \end{gathered}
\end{equation}

\begin{theorem}
(updated version of Proposition 5.2) $H^1$ energy estimates. Can we have an analogous structure when multiplying the above equations by the Laplacians?
\end{theorem}

\begin{proof}

We know already that the original (coloured black) parts of (\ref{U_eps_C_eps}), multiplied by $\Delta V^\epsilon, \Delta W^\epsilon$ yield

\begin{equation} \label{originalH1}
  \begin{split}
    & \frac{1}{2} \frac{\diff}{\diff t} (\norm{\nabla V^\epsilon}^2_2 + \norm{\epsilon \nabla W^\epsilon}^2_2) + \frac{3}{5} (\norm{\Delta V^\epsilon}^2_2 + \norm{\epsilon \Delta W^\epsilon}^2_2) \\
    & \leq C_1 (\norm{\nabla V^\epsilon}^2_2 + \norm{\epsilon \nabla W^\epsilon}^2_2) \big [ (\norm{\Delta V}^2_2 + \norm{\epsilon \Delta W^\epsilon}^2_2) + (1+ \epsilon^4) \norm{\Delta v}^2_2 \norm{\nabla \Delta v}^2_2 \big ] \\
    & + C_1 \epsilon^2 (1 + \norm{\Delta v}^2_2) (\norm{\nabla \partial_t v}^2_2 + \norm{\nabla \Delta v}^2_2).
  \end{split}
\end{equation}

Now we examine the effect of the new contributions from the concentration equation and the additional Coriolis terms.

\textcolor{green}{\textbf{Part one: the concentration function.}}

Multiply the concentration equation in (\ref{U_eps_C_eps}) by $- \Delta C^\epsilon$ and integrate by parts.

On the left-hand side we have \textit{i)} the contribution from $\partial_t C^\epsilon \cdot \Delta C^\epsilon,$ which is $\frac{1}{2}\frac{\diff}{\diff t} \norm{\nabla C^\epsilon}^2_2$ and \textit{ii)} the contribution from the diffusivity terms.

$$ \frac{1}{2}\frac{\diff}{\diff t} \norm{\nabla C^\epsilon}^2_2 + K \norm{\Delta C^\epsilon}_2^2 \leq \int_\Omega \big ( U^\epsilon \nabla C^\epsilon  + U^\epsilon \nabla c+ u \nabla C^\epsilon \big ) \Delta C^\epsilon - \int_\Omega S^\epsilon \Delta C^\epsilon$$

To estimate $\int_\Omega \big ( U^\epsilon \nabla C^\epsilon  + U^\epsilon \nabla c + u \nabla C^\epsilon \big ) \Delta C^\epsilon$ we use Lemma 2.2. \todo{why does it require zero spatial average even for the case of periodic domain}
We use the generalised Young inequality with a constant $\zeta$ that is necessary so that later we can bring certain terms to the left-hand side.

Firstly, we have

\begin{equation}
  \begin{split}
    & \int_\Omega U^\epsilon \nabla C^\epsilon \Delta C^\epsilon \leq \norm{\nabla V^\epsilon}_2^{1/2} \norm{\Delta V^\epsilon}_2^{1/2} \norm{\nabla C^\epsilon}_2^{1/2} \norm{\Delta C^\epsilon}_2^{1/2} \norm{\Delta C^\epsilon}_2 \\
    & \leq \zeta \norm{\Delta C^\epsilon}_2^2 + \frac{1}{\zeta} \norm{\nabla V^\epsilon}_2 \norm{\Delta V^\epsilon}_2 \norm{\nabla C^\epsilon}_2 \norm{\Delta C^\epsilon}_2 \\
    & \leq \zeta \norm{\Delta C^\epsilon}_2^2 + \frac{1}{\zeta} \big ( \frac{1}{2} \norm{\nabla V^\epsilon}^2_2 \norm{\Delta V^\epsilon}^2_2 + \frac{1}{2} \norm{\nabla C^\epsilon}^2_2 \norm{\Delta C^\epsilon}^2_2 \big ).
  \end{split}
\end{equation}

Secondly,

\begin{equation}
  \begin{split}
    & \int_\Omega U^\epsilon \nabla c \Delta C^\epsilon \leq \norm{\nabla V^\epsilon}_2^{1/2} \norm{\Delta V^\epsilon}_2^{1/2} \norm{\nabla c}_2^{1/2} \norm{\Delta c}_2^{1/2} \norm{\Delta C^\epsilon}_2 \\
    & \zeta_1 \norm{\Delta C^\epsilon}^2_2 + \frac{1}{\zeta_1} \norm{\nabla V^\epsilon}_2 \norm{\Delta V^\epsilon}_2 \norm{\nabla c}_2 \norm{\Delta c}_2 \\
    & \zeta_1 \norm{\Delta C^\epsilon}^2_2 + \frac{1}{\zeta_1} \big (\zeta_2 \norm{\Delta V^\epsilon}_2^2 + \frac{1}{\zeta_2}  \norm{\nabla V^\epsilon}^2_2  \norm{\nabla c}^2_2 \norm{\Delta c}^2_2 \big ).
  \end{split}
\end{equation}

Finally here we use the Young inequality with $\frac{1}{4} + \frac{3}{4},$ since we have $\norm{\Delta C^\epsilon}_2^{1/2} \norm{\Delta C^\epsilon}_2 = \norm{\Delta C^\epsilon}_2^{3/2}.$

\begin{equation}
  \begin{split}
    & \int_\Omega u \nabla C^\epsilon \Delta C^\epsilon \leq \norm{\nabla v}_2^{1/2} \norm{\Delta v}_2^{1/2} \norm{\nabla C^\epsilon}_2^{1/2} \norm{\Delta C^\epsilon}_2^{1/2} \norm{\Delta C^\epsilon}_2 \\
    & \leq \zeta \norm{\Delta C^\epsilon}_2^2 + \frac{1}{\zeta} \norm{\nabla v}^2_2 \norm{\Delta v}^2_2 \norm{\nabla C^\epsilon}^2_2 \\
    & \leq \zeta \norm{\Delta C^\epsilon}_2^2 + \frac{1}{\zeta} \norm{\Delta v}^2_2 \norm{\nabla \Delta v}^2_2 \norm{\nabla C^\epsilon}^2_2.
  \end{split}
\end{equation}

Source \todo{todo} term.

With this we arrive to first updated version of the inequality (\ref{originalH1}), which now becomes

\begin{equation}
  \begin{split}
    & \frac{1}{2} \frac{\diff}{\diff t} (\norm{\nabla V^\epsilon}^2_2 + \norm{\epsilon \nabla W^\epsilon}^2_2 + \textcolor{green}{\norm{\nabla C^\epsilon}^2_2}) + \frac{3}{5} (\norm{\Delta V^\epsilon}^2_2 + \norm{\epsilon \Delta W^\epsilon}^2_2 \textcolor{green}{+ K \norm{\Delta C^\epsilon}_2^2}) \\
    & \leq C_1 (\norm{\nabla V^\epsilon}^2_2 + \norm{\epsilon \nabla W^\epsilon}^2_2 \textcolor{green}{+ \norm{\nabla C^\epsilon}^2_2}) \\
    & \times \big [ (\norm{\Delta V}^2_2 + \norm{\epsilon \Delta W^\epsilon}^2_2 \textcolor{green}{+ K \norm{\Delta C^\epsilon}_2^2}) + (1+ \epsilon^4) \norm{\Delta v}^2_2 \norm{\nabla \Delta v}^2_2 + \textcolor{green}{\norm{\nabla c}^2_2 \norm{\Delta c}^2_2} \big ] \\
    & + C_1 \epsilon^2 (1 + \norm{\Delta v}^2_2) (\norm{\nabla \partial_t v}^2_2 + \norm{\nabla \Delta v}^2_2),
  \end{split}
\end{equation}

after taking the $\zeta \cdot \text{Laplacian}$- type terms such as $\zeta \norm{\Delta C^\epsilon}_2^2, \frac{1}{\zeta_1} \zeta_2 \norm{\Delta V^\epsilon}_2^2$ to the left-hand side.

This is the \todo{todo: $\norm{\Delta c}^2_2$ boundedness} version of the inequality after adding the pollution equation effects, without the source term.

\textcolor{blue}{\textbf{Part two: the Coriolis term.}}

Here we calculate the contributions from the Coriolis terms. In this part we take advantage of the Young inequality and being able to pass from $w, W^\epsilon$ in norms to $\nabla v, \nabla V^\epsilon$ because of the divergence-free condition.

\begin{enumerate}

  \item $$\int_\Omega \gamma V^\epsilon \cdot \Delta V^\epsilon \leq C \norm{\nabla V^\epsilon}^2_2$$
  
  \item $$ \int_\Omega \epsilon \beta W^\epsilon \Delta V^\epsilon \leq C \norm{\epsilon \nabla W^\epsilon}_2 \norm{\nabla V^\epsilon}_2 \leq C \norm{\epsilon \nabla W^\epsilon}^2_2 + C \norm{\nabla V^\epsilon}^2_2.$$
  
  \item $$ \int_\Omega \epsilon \beta w \Delta V^\epsilon \leq C \epsilon \norm{w}_2 \cdot \norm{\Delta V^\epsilon}_2 \leq C \cdot \epsilon \norm{\nabla v}_2 \norm{\Delta V^\epsilon}_2 $$
  $$ \leq C \frac{1}{\zeta} \epsilon^2 \norm{\nabla v}_2^2 + \zeta \norm{\Delta V^\epsilon}^2_2.$$
  
  \item $$ \int_\Omega \epsilon \beta V^\epsilon_1 \Delta W^\epsilon \leq C \norm{\nabla V^\epsilon_1}_2 \norm{\epsilon \nabla W^\epsilon}_2 \leq C \norm{\nabla V^\epsilon_1}^2_2 + C \norm{\epsilon \nabla W^\epsilon}^2_2.$$
  
  \item $$ \int_\Omega \epsilon \beta v_1 \Delta W^\epsilon \leq C \norm{\epsilon \nabla v_1}_2 \norm{\nabla W^\epsilon}_2 \leq C \norm{\epsilon \nabla v_1}_2 \norm{\Delta V^\epsilon}_2 $$
  $$ \leq C \frac{1}{\zeta} \epsilon^2 \norm{\nabla v}_2^2 + \zeta \norm{\Delta V^\epsilon}^2_2.$$
  
\end{enumerate}

Regarding these inequalities and integrating them into the the main one, note that \textit{a)} the standalone Laplacians can go to the left-hand side, \textit{b)}  the $C \cdot \norm{\nabla V^\epsilon}^2_2$ -type terms can be added to the part where $\norm{\nabla V^\epsilon}^2_2$ is already present, and \textit{c)} the $\epsilon^2 \norm{\nabla v}^2_2$ can be added to the second part where $\epsilon^2$ is a multiplier.

Connecting all this, we have the updated version of the inequality which was the inequality after (5.7) in Li-Titi. Now it takes the following form:

\begin{equation}
  \begin{split}
    & \frac{1}{2} \frac{\diff}{\diff t} (\norm{\nabla V^\epsilon}^2_2 + \norm{\epsilon \nabla W^\epsilon}^2_2 + \textcolor{green}{\norm{\nabla C^\epsilon}^2_2}) + \frac{3}{5} (\norm{\Delta V^\epsilon}^2_2 + \norm{\epsilon \Delta W^\epsilon}^2_2 \textcolor{green}{+ K \norm{\Delta C^\epsilon}_2^2}) \\
    & \leq C_1 (\norm{\nabla V^\epsilon}^2_2 + \norm{\epsilon \nabla W^\epsilon}^2_2 \textcolor{green}{+ \norm{\nabla C^\epsilon}^2_2}) \\
    & \times \big [ (\norm{\Delta V}^2_2 + \norm{\epsilon \Delta W^\epsilon}^2_2 \textcolor{green}{+ K \norm{\Delta C^\epsilon}_2^2}) + \textcolor{blue}{C}+ (1+ \epsilon^4) \norm{\Delta v}^2_2 \norm{\nabla \Delta v}^2_2 + \textcolor{green}{\norm{\nabla c}^2_2 \norm{\Delta c}^2_2} \big ] \\
    & + C_1 \epsilon^2 (1 + \norm{\Delta v}^2_2) (\norm{\nabla \partial_t v}^2_2 + \norm{\nabla \Delta v}^2_2 + \textcolor{blue}{\norm{\nabla v}_2^2}).
  \end{split}
\end{equation}

As long as $\norm{\nabla V^\epsilon}^2_2 + \norm{\epsilon \nabla W^\epsilon}^2_2 + \textcolor{green}{\norm{\nabla C^\epsilon}^2_2}) \leq \delta_0^2,$ with $\delta_0 = \sqrt{\frac{1}{10 C_1}}$ then 

$$ C_1 (\norm{\nabla V^\epsilon}^2_2 + \norm{\epsilon \nabla W^\epsilon}^2_2 \textcolor{green}{+ \norm{\nabla C^\epsilon}^2_2})(\norm{\Delta V}^2_2 + \norm{\epsilon \Delta W^\epsilon}^2_2 \textcolor{green}{+ K \norm{\Delta C^\epsilon}_2^2}) $$

$$ \leq \delta_0^2 (\norm{\Delta V}^2_2 + \norm{\epsilon \Delta W^\epsilon}^2_2 \textcolor{green}{+ K \norm{\Delta C^\epsilon}_2^2}) $$

can be "melted into" the left-hand side and we finally have

\begin{equation}
  \begin{split}
    & \frac{\diff}{\diff t} (\norm{\nabla V^\epsilon}^2_2 + \norm{\epsilon \nabla W^\epsilon}^2_2 + \textcolor{green}{\norm{\nabla C^\epsilon}^2_2}) + (\norm{\Delta V^\epsilon}^2_2 + \norm{\epsilon \Delta W^\epsilon}^2_2 \textcolor{green}{+ K \norm{\Delta C^\epsilon}_2^2}) \\
    & \leq C_1 \big [ \textcolor{blue}{C}+ (1+ \epsilon^4) \norm{\Delta v}^2_2 \norm{\nabla \Delta v}^2_2 + \textcolor{green}{\norm{\nabla c}^2_2 \norm{\Delta c}^2_2} \big ] \\
    & \times (\norm{\nabla V^\epsilon}^2_2 + \norm{\epsilon \nabla W^\epsilon}^2_2 \textcolor{green}{+ \norm{\nabla C^\epsilon}^2_2}) \\
    & + C_1 \epsilon^2 (1 + \norm{\Delta v}^2_2) (\norm{\nabla \partial_t v}^2_2 + \norm{\nabla \Delta v}^2_2 + \textcolor{blue}{\norm{\nabla v}_2^2}).
  \end{split}
\end{equation}

Using $(V^\epsilon, W^\epsilon, C^\epsilon)(t_0) = 0,$ applying the Gronwall inequality we have 

\begin{equation} \label{result_theorem1}
  \begin{split}
    & \sup_{0 \leq s \leq t}{(\norm{\nabla V^\epsilon}^2_2 + \norm{\epsilon \nabla W^\epsilon}^2_2 + \textcolor{green}{\norm{\nabla C^\epsilon}^2_2}) } + \int_0^t (\norm{\Delta V^\epsilon}^2_2 + \norm{\epsilon \Delta W^\epsilon}^2_2 \textcolor{green}{+ K \norm{\Delta C^\epsilon}^2_2}) \diff s \\
    & \leq C_1 \epsilon^2 \bigg ( \int_0^t (1 + \norm{\Delta v}^2_2) (\norm{\nabla \partial_t v}^2_2 + \norm{\nabla \Delta v}^2_2 + \textcolor{blue}{\norm{\nabla v}_2^2}) \diff s \bigg ) \\
    & \times \exp{ \int_0^t \textcolor{blue}{C}+ (1+ \epsilon^4) \norm{\Delta v}^2_2 \norm{\nabla \Delta v}^2_2 + \textcolor{green}{\norm{\nabla c}^2_2 \norm{\Delta c}^2_2} \diff s}.
  \end{split}
\end{equation}

\end{proof}

Now we move on to check if the extended, analogous theorem to Proposition 5.3 works.

\begin{theorem}
When $\epsilon$ is small enough, i.e. $\epsilon \in (0, \epsilon_0),$ then we have global unique strong solutions for the extended system, and we have

$$ \sup_{0 \leq t < \infty}{(\norm{V^\epsilon}^2_{H^1} + \norm{\epsilon W^\epsilon}^2_{H^1} + \textcolor{green}{\norm{C^\epsilon}^2_{H^1}}) } + \int_0^\infty (\norm{\nabla V^\epsilon}^2_{H^1} + \norm{\epsilon \nabla W^\epsilon}^2_{H^1} \textcolor{green}{+ K \norm{\nabla C^\epsilon}^2_{H^1}}) \diff t \leq C \epsilon^2. $$

\end{theorem}

\begin{proof}

This part is really analogous to what is written on p27, \textbf{the main todo parts are}:

\begin{itemize}

  \item POINT A) We need the concentration-including version of Corollary 3.1, namely, an estimate for the term $\int_0^t \norm{\nabla c}^2_2 \norm{\Delta c}^2_2$ in the integral in the above result (\ref{result_theorem1}). Corollary 3.1 gives a constant estimate only for the velocities, e.g. $\int_0^\infty \norm{\nabla v}^2_{H^2}.$
  \item POINT B) We need the $L^2$ estimate (analogous to Proposition 5.1.) to $C^\epsilon,$ while we already know the result for $V^\epsilon, \epsilon W^\epsilon.$
  \item POINT C) We need an existence result for (E-SNS) that states the unique local existence of strong solution $(u^\epsilon, c^\epsilon)$.
  \item POINT D) We need an energy-existence result for (E-SNS) that describes the connection between maximal existence time and the blow-up of the energy of the system. In other words, a result that states that at a certain time $T^*$ the energy is small, then the solution can be continued.

\end{itemize}

When these missing pieces are clear, then from the $L^2$ estimate (POINT B) we have

$$ \sup_{0 \leq t < T^*_\epsilon}{(\norm{V^\epsilon}^2_2 + \norm{\epsilon W^\epsilon}^2_2 + \textcolor{green}{\norm{C^\epsilon}^2_2}) } + \int_0^{T^*_\epsilon} (\norm{\nabla V^\epsilon}^2_2 + \norm{\epsilon \nabla W^\epsilon}^2_2 \textcolor{green}{+ K \norm{\nabla C^\epsilon}^2_2}) \diff t \leq C \epsilon^2, $$

$T^*_\epsilon$ being the maximal existence time of $(u^\epsilon, c^\epsilon)$ (POINT C).

From this, using the previously defined $\delta_0,$ we define

$$t^*_\epsilon:= \sup{ \big \{ t \in (0, T^*_\epsilon) \mid \sup_{0 \leq s \leq t}{\big (\norm{\nabla V^\epsilon}^2_2 + \norm{\epsilon \nabla W^\epsilon}^2_2 \textcolor{green}{+ K \norm{\nabla C^\epsilon}^2_2} \big ) \leq \delta_0^2}\} }$$

Using POINT A), from (\ref{result_theorem1}) we get

\begin{equation}
  \begin{split}
    & \sup_{0 \leq s \leq t}{(\norm{\nabla V^\epsilon}^2_2 + \norm{\epsilon \nabla W^\epsilon}^2_2 + \textcolor{green}{\norm{\nabla C^\epsilon}^2_2}) } \\
    & + \int_0^t (\norm{\Delta V^\epsilon}^2_2 + \norm{\epsilon \Delta W^\epsilon}^2_2 \textcolor{green}{+ K \norm{\Delta C^\epsilon}^2_2}) \diff s \leq C \epsilon^2
  \end{split}
\end{equation}

for any $t \in [0, t^*_\epsilon).$ This is important: in order to apply (\ref{result_theorem1}), we need to make sure that its condition holds, i.e. that we are in $[0, t^*_\epsilon),$ which is the range where it applies.

Setting $\epsilon_0 = \frac{\delta_0}{\sqrt{2C}},$ we have

\begin{equation}
  \begin{split}
    & \sup_{0 \leq s \leq t}{(\norm{\nabla V^\epsilon}^2_2 + \norm{\epsilon \nabla W^\epsilon}^2_2 + \textcolor{green}{\norm{\nabla C^\epsilon}^2_2}) } \\
    & + \int_0^t (\norm{\Delta V^\epsilon}^2_2 + \norm{\epsilon \Delta W^\epsilon}^2_2 \textcolor{green}{+ K \norm{\Delta C^\epsilon}^2_2}) \diff s \leq \frac{\delta_0}{2}
  \end{split}
\end{equation}

for any $\epsilon \in (0, \epsilon_0)$ and any $t \in [0, t^*_\epsilon).$ From which we have

$$ \sup_{0 \leq s < t^*_\epsilon}{(\norm{\nabla V^\epsilon}^2_2 + \norm{\epsilon \nabla W^\epsilon}^2_2 + \textcolor{green}{\norm{\nabla C^\epsilon}^2_2}) } \leq \frac{\delta_0}{2}.$$

From this we have $t^*_\epsilon = T^*_\epsilon,$ and as a consequence,

\begin{equation}
  \begin{split}
    & \sup_{0 \leq s < T^*_\epsilon}{(\norm{\nabla V^\epsilon}^2_2 + \norm{\epsilon \nabla W^\epsilon}^2_2 + \textcolor{green}{\norm{\nabla C^\epsilon}^2_2}) } \\
    & + \int_0^{T^*_\epsilon} (\norm{\Delta V^\epsilon}^2_2 + \norm{\epsilon \Delta W^\epsilon}^2_2 \textcolor{green}{+ K \norm{\Delta C^\epsilon}^2_2}) \diff s \leq C \epsilon^2.
  \end{split}
\end{equation}

From this we have $T^*_\epsilon = \infty$ (using POINT D).

\end{proof}

\subsection{Questions and todo}

\begin{itemize}

  \item Initial and boundary conditions, for example p27, they use (1.2)-(1.4). Again, are these symmetries important (the even and odd functions?) Plus, define BC, IC for the $c^\epsilon, c$ functions.
  
  \item How exactly do we know that $\big (\norm{\nabla V^\epsilon}^2_2 + \norm{\epsilon \nabla W^\epsilon}^2_2 \textcolor{green}{+ K \norm{\nabla C^\epsilon}^2_2} \big )$ grows continuously? If it didn't, then it would be possible that from $\delta_0 / 2$ it jumps to $\delta_0 +1,$ and with that we wouldn't have $t^*_\epsilon = T^*_\epsilon.$
  
  \item In Proposition 5.3 why exactly do we need the $L^2$ result from Proposition 5.1? Is it because for the bound of $\int_0^\infty \norm{\nabla V^\epsilon}^2_{H^1}$ by definition we need $\int_0^\infty \norm{\nabla V^\epsilon}^2_{L^2}$ and $\int_0^\infty \norm{\Delta V^\epsilon}^2_{L^2}$, and the first does not come from the second one by the Poincare inequality?
  
  \item on p27 why do we need (5.8), which is essentially Proposition 5.1 in order to have $t^*_\epsilon = T^*_\epsilon?$ If it is only needed towards the end, why is it mentioned in the beginning?

\end{itemize}

\section{Downwind-matching coordinates}

(E-SNS): extended scaled Navier-Stokes.

    \begin{equation} \label{ANSC1}
    \partial_t u^\epsilon _1 + \mathbf{u}^\epsilon \cdot \nabla u^\epsilon _1 -\Delta_\nu u^\epsilon _1 -\gamma u_2^\epsilon + \epsilon \beta u_3^\epsilon + \partial_1 p^\epsilon = 0 \text{ in } \Omega \times (0,T)
    \end{equation}
    
    \begin{equation} \label{ANSC2}
    \partial_t u_2^\epsilon + \mathbf{u}^\epsilon \cdot \nabla u_2^\epsilon -\Delta_\nu u_2^\epsilon + \gamma u_1^\epsilon  - \alpha \epsilon u^\epsilon_3 + \partial_2 p^\epsilon = 0 \text{ in } \Omega \times (0,T)
    \end{equation}
    
    \begin{equation} \label{ANSC3}
    \epsilon^2 \{  \partial_t u_3^\epsilon + \mathbf{u}^\epsilon \cdot \nabla u_3^\epsilon -\Delta_\nu u_3^{\epsilon} \} + \epsilon \alpha u^\epsilon_2 - \epsilon \beta u_1^\epsilon + \partial_3 p^\epsilon = 0 \text{ in } \Omega \times (0,T)
    \end{equation}
    
    \begin{equation} \label{ANSC4}
    \nabla \cdot \mathbf{u}^\epsilon= 0 \text{ in } \Omega \times (0,T)
    \end{equation}
    
    \begin{equation} \label{ANSC5}
    \partial_t c^\epsilon + \mathbf{u^\epsilon} \cdot \nabla c^\epsilon = \epsilon K_1 \partial_{11} ^ 2 c^\epsilon + K_2 \partial_{22} ^ 2 c^\epsilon + K_3 \partial_{33} ^ 2 c^\epsilon + s^\epsilon \text{ in } \Omega \times (0,T).
    \end{equation}
    

(E-PEs): extended Primitive Equations.

    \begin{equation} \label{HNSC1}    
    \partial_t u_1 + \mathbf{u} \cdot \nabla u_1 - \Delta_\nu u_1 - \gamma u_2 + \partial_1 p = 0 \text{ in } \Omega \times (0,T)
    \end{equation}
    
    \begin{equation} \label{HNSC2}   
    \partial_t u_2 + \mathbf{u} \cdot \nabla u_2 - \Delta_\nu u_2 + \gamma u_1 + \partial_2 p = 0 \text{ in } \Omega \times (0,T)
    \end{equation}
    
    \begin{equation} \label{HNSC3}   
    \partial_3 p = 0 \text{ in } \Omega \times (0,T)
    \end{equation}
    
    \begin{equation} \label{HNSC4}   
    \nabla \cdot \mathbf{u} = 0 \text{ in } \Omega \times (0,T)
    \end{equation}
    
    \begin{equation} \label{HNSC5}   
    \partial_t c + \mathbf{u} \cdot \nabla c = K_2 \partial_{22} ^ 2 c + K_3 \partial_{33} ^2 c + s \text{ in } \Omega \times (0,T)
    \end{equation}
    
The main difference in this case is that in the equations that govern $(U^\epsilon, C^\epsilon),$ we have

$$ \partial_t C^\epsilon + U^\epsilon \nabla C^\epsilon  + U^\epsilon \nabla c+ u \nabla C^\epsilon = \textcolor{orange}{\epsilon K_1 \partial_{11} C^\epsilon + \epsilon K_1 \partial_{11} c} + K_2 \partial_{22} C^\epsilon + K_3 \partial_{33} C^\epsilon + S^\epsilon. $$

This means that we can not follow the same path as before, as multiplying the equation by $- \Delta C^\epsilon$ results only $\epsilon \norm{\partial_{11} C^\epsilon}$ on the left-hand side, while from the estimates of the convective terms, we still have $\zeta \norm{\Delta V^\epsilon}^2_2$ terms on the right-hand side of the estimate, and these can not anymore be "melted" into the left-hand side, as the full, $\epsilon$-free Laplacian is missing.


\end{document}

